\section{The formal verification}\label{app:formal}

The formalization uses Lean~4~\cite{Lean4} and Mathlib~\cite{Mathlib20} and is in the directory
\lean{Morrey/HomogeneousIsotropic} of the repository~\cite{Companion}. We describe its
statements, their verification and dependencies, the points where they differ from the text, and
the correspondence between the paper and the Lean declarations.

\subsection{The statements}\label{app:formal-statements}

{\sloppy
The file \lean{palomar/HomogeneousIsotropic/Challenge.lean} imports only Mathlib and states the
results of \cref{sec:intro} as theorems of the namespace \lean{HomogeneousIsotropicChallenge}.
They are
\lean{theorem1_1}, \lean{theorem1_2}, \lean{theorem1_3_a}, \lean{theorem1_3_b},
\lean{corollary1_4}, \lean{theorem1_5}, \lean{theorem1_5_identity}, \lean{corollary1_6_a},
\lean{corollary1_6_b}, \lean{proposition1_7} and \lean{corollary1_8}, together with the
identity $\det A=\abs z^2-\abs w^2$ (\lean{det_eq_norm_sq_sub}), which ties the coordinates to
the determinant. Its definitions are the following.\par}
\begin{itemize}
\item Matrices are \leanx{Matrix (Fin 2) (Fin 2) $\R$}, and the gradient of
  $\varphi\colon\R^2\to\R^2$ is the matrix \leanx{jac $\varphi$ x} of the derivatives of its
  components.
\item \leanx{IsRankOneConvex F}: $\tau\mapsto F(A+\tau B)$ is convex on $\R$ for every $A$ and
  every $B$ with \leanx{B.rank~=~1}.
\item \leanx{IsQuasiconvex F}: $F(A)\le\int_{(0,1)^2}F(A+\nabla\varphi)$ for every $A$ and every
  Lipschitz $\varphi$ vanishing outside the open unit square, for which $F(A+\nabla\varphi)$ is
  integrable.
\item \leanx{IsIsotropic F}: invariance under
  \leanx{Matrix.specialOrthogonalGroup (Fin 2) $\R$} on both sides. \leanx{IsPosHomogeneous p F}:
  $F(cA)=c^pF(A)$ for $c>0$, with the real power.
\item \leanx{Polyconvex F}: $F(A)=g(A,\det A)$ for a convex $g$ on the product of the matrices
  and $\R$. \leanx{HasPolyaffineSupportAt F A$_0$}: a function
  $k+\sum_{i,j}K_{ij}A_{ij}+c\det A$ lies below $F$ and equals $F$ at $A_0$.
\item \lean{conformalPart}, \lean{anticonformalPart}: the coordinates $z$, $w$.
  \leanx{burkholderType p c$_0$ m} is $G_{c_0,m}$, written with $\abs z+\abs w$ and
  $\abs z-\abs w$, and \leanx{burkholder p A} is
  $(\abs z+\abs w)^{p-1}(\abs z-(p-1)\abs w)=\tfrac p2B_p(A)$ for $1\le p\le2$.
\item \leanx{multiplier p H m t} is the number $c^\ast$ of \cref{thm:closed}.
\end{itemize}
The definitions are copied verbatim from the library, so the library's theorems have precisely
the statements given in the Challenge file.

\subsection{Verification and dependencies}\label{app:formal-check}

{\sloppy
The theorems are proved in \lean{palomar/HomogeneousIsotropic/Solution.lean} from the library.
Lean's kernel checks the proofs. Every theorem depends only on the standard axioms \lean{propext},
\lean{Classical.choice} and \lean{Quot.sound}, as verified by an audit of every declaration of
the library. Lean's Comparator, with the configuration
\lean{palomar/HomogeneousIsotropic/comparator.json}, checks that the Solution proves the
statements of the Challenge with the same definitions, under Lean's kernel and the independent kernel nanoda. Lean
and Mathlib are pinned at version \texttt{v4.36.0-rc2}.\par}

The library's proofs use the rank-one convexity and the quasiconvexity of the library
\lean{CalculusOfVariations}~\cite{CoVLean}, which are the notions above
(\lean{Morrey/HomogeneousIsotropic/Bridges.lean}). From that library they take two facts: minors
are null Lagrangians, which gives \cref{lem:supports}\,(a), and rank-one convex functions are
locally Lipschitz, which gives the continuity used in \cref{prop:small}. From the
formalization~\cite{IwaniecLean} of~\cite{SharpHI} they take one theorem,
\lean{Agazzi.quasiconvex_burkholder}: for $1<p<2$ the function
$(\abs z+\abs w)^{p-1}(\abs z-(p-1)\abs w)$ is quasiconvex in the sense of
\lean{CalculusOfVariations}. It is used in \lean{Morrey/HomogeneousIsotropic/Burkholder.lean},
for \cref{cor:low}\,(a) and \cref{cor:all}. \Cref{lem:convex,lem:sublinear}
are in the file \lean{Morrey/Common/Basic/Analysis/Convex/Tangent.lean}. The hypotheses
of \cref{thm:main} are checked on an integrand that is not convex, $F(A)=(\det A)^2$ with $p=4$
(\lean{polyconvex_detSq}, \lean{not_convexOn_detSq}).

\subsection{Differences from the text}\label{app:formal-differences}

\begin{itemize}
\item The formal statements are written in conformal coordinates: they have $\abs z+\abs w$
  where the text has $\abs A$. The identity $\det A=\abs z^2-\abs w^2$ is formalized; the
  identity $\abs A=\abs z+\abs w$ for the operator norm is not.
\item Of \cref{lem:coordinates}\,(b) and of the description of isotropic functions after it, the
  library has the implications that the proofs use: a nonzero matrix with
  $\abs\alpha=\abs\beta$ has rank one, and an isotropic $F$ is $\Phi(\abs z,\abs w)$.
\item The formalization~\cite{IwaniecLean} proves the quasiconvexity of the Burkholder function
  for Lipschitz periodic test maps, by the arguments of~\cite{SharpHI}, and passes to the test
  maps used here by the equivalence of the two notions, which \lean{CalculusOfVariations}
  proves (\cref{sec:intro-low}). The identity of \cref{thm:low} is stated with
  \leanx{burkholder p}, hence with the denominators $p(2-p)$.
\item In \cref{thm:closed} the supporting line is given by its slope: the hypothesis is
  $h(s)\ge H+m(s-t)$ for $s\in[-1,1]$.
\item The equalities at $(x_t,y_t)$ in \cref{lem:minorant}, in \cref{prop:caseI}, in Case II and
  in \cref{prop:clamp} are not part of the declarations listed for them. They are proved where
  they are used, in \lean{IsProfile.exists_support_of_add_eq_one} and in
  \lean{polyaffine_multiplier_le}. The support at $t=-1$ is obtained from the one at $t=1$
  inside the same proof.
\item Some formal statements are more general: the identity of \cref{thm:low} is stated for
  $0<p<2$, and \cref{prop:tangent}, \cref{cor:BE} and \cref{lem:ends} for every real $p$. The
  library also has the implication behind \cref{cor:low}\,(a), for $1\le p<2$: if the Burkholder
  function is quasiconvex, so is $F$ (\lean{quasiconvex_of_rankOneConvex_of_burkholder}).
\item The facts recalled in \cref{sec:intro}, such as the rank-one convexity of $B_p$ and the
  formula $F(A)=\abs A^p\,h(\det A/\abs A^2)$, are not formalized.
\end{itemize}

\subsection{Dictionary}\label{app:formal-dictionary}

\Cref{tab:formal-dictionary} lists the Lean declarations corresponding to the numbered results
of the paper. Names are relative to the namespace \lean{HomogeneousIsotropic}; the prefix
\lean{HomogeneousIsotropicChallenge} marks the statements of the Challenge file, and the two
declarations with the prefix \lean{ConvexOn} are in the root namespace. Multiple declarations in one row
jointly establish the corresponding result.

{\small
\begin{longtable}{@{}>{\raggedright\arraybackslash}p{0.25\textwidth}>{\raggedright\arraybackslash}p{0.71\textwidth}@{}}
\caption{The Lean declarations behind the numbered results.}\label{tab:formal-dictionary}\\
\toprule
Result & Lean declarations\\
\midrule
\endfirsthead
\toprule
Result & Lean declarations\\
\midrule
\endhead
\bottomrule
\endfoot
\multicolumn{2}{@{}l}{\emph{\Cref{sec:intro}}}\\*[2pt]
\Cref{thm:main} & \lean{HomogeneousIsotropicChallenge.theorem1_1}; \lean{hasPolyaffineSupportAt_of_rankOneConvex}, \lean{polyconvex_of_rankOneConvex}, \lean{quasiconvex_of_rankOneConvex}\\
\Cref{thm:closed} & \lean{HomogeneousIsotropicChallenge.theorem1_2}; \lean{polyaffine_multiplier_le}\\
\Cref{thm:iff} & \lean{HomogeneousIsotropicChallenge.theorem1_3_a}, \lean{HomogeneousIsotropicChallenge.theorem1_3_b}; \lean{polyconvex_iff_of_two_lt}, \lean{nonneg_of_polyconvex}, \lean{polyconvex_iff_two}\\
\Cref{cor:equiv} & \lean{HomogeneousIsotropicChallenge.corollary1_4}; \lean{tfae_of_two_le}\\
\Cref{thm:low} & \lean{HomogeneousIsotropicChallenge.theorem1_5}, \lean{HomogeneousIsotropicChallenge.theorem1_5_identity}; \lean{exists_burkholderType_le}, \lean{burkholderType_eq}\\
\Cref{cor:low}\,(a) & \lean{HomogeneousIsotropicChallenge.corollary1_6_a}; \lean{quasiconvex_burkholder}, \lean{quasiconvex_of_rankOneConvex_of_lt_two}\\
\Cref{cor:low}\,(b) & \lean{HomogeneousIsotropicChallenge.corollary1_6_b}; \lean{nonneg_of_rankOneConvex_one}, \lean{convexOn_of_rankOneConvex_one}\\
\Cref{prop:small} & \lean{HomogeneousIsotropicChallenge.proposition1_7}; \lean{eq_zero_of_rankOneConvex_of_lt_one}\\
\Cref{cor:all} & \lean{HomogeneousIsotropicChallenge.corollary1_8}; \lean{quasiconvex_of_rankOneConvex_of_pos}\\[4pt]
\multicolumn{2}{@{}l}{\emph{\Cref{sec:setting,app:setting}}}\\*[2pt]
\Cref{lem:coordinates} & \lean{HomogeneousIsotropicChallenge.det_eq_norm_sq_sub}; \lean{det_ofComplex}, \lean{rank_ofComplex_eq_one}, \lean{ofComplex_mul_mul}, \lean{conformalPart_reflection_mul}\\
\Cref{lem:profile} & \lean{IsIsotropic.eq_profile}, \lean{isProfile_profile}\\
\Cref{lem:supports} & \lean{HasPolyaffineSupportAt.quasiconvexAtOn}, \lean{polyconvex_of_hasPolyaffineSupportAt}, \lean{Polyconvex.hasPolyaffineSupportAt}, \lean{Polyconvex.quasiconvex}, \lean{Polyconvex.rankOneConvex}\\
\Cref{lem:lift} & \lean{hasPolyaffineSupportAt_of_profile}\\[4pt]
\multicolumn{2}{@{}l}{\emph{\Cref{sec:roc,app:roc}}}\\*[2pt]
\Cref{lem:basic} & \lean{IsProfile.convexOn_slice}, \lean{IsProfile.le_of_sq_sub_sq_eq}\\
\Cref{lem:minorant} & \lean{IsProfile.exists_burkholderMinorant}, \lean{IsProfile.burkholderMinorant_of_slice}\\
\Cref{lem:convex} & \lean{ConvexOn.add_mul_le_of_hasDerivAt_le}\\
\Cref{prop:tangent} & \lean{IsProfile.tangent}\\
\Cref{cor:BE} & \lean{IsProfile.baker_ericksen}\\
\Cref{lem:ends} & \lean{IsProfile.slice_ge_of_one}, \lean{IsProfile.slice_ge_of_neg_one}\\[4pt]
\multicolumn{2}{@{}l}{\emph{\Cref{sec:pge2,app:pge2}}}\\*[2pt]
\Cref{lem:burkholder} & \lean{burkholder_support_le}\\
\eqref{eq:caseIII} & \lean{BurkholderMinorant.support}, \lean{BurkholderMinorant.support_two}\\
\Cref{lem:caseI-algebra} & \lean{caseOne_slope_pos}, \lean{caseOne_const_neg}, \lean{caseOne_relations}\\
\Cref{lem:mu} & \lean{caseOne_slope_bounds}\\
\Cref{lem:line} & \lean{IsProfile.support_of_stationary}\\
\Cref{lem:psi} & \lean{exists_stationary_ge}\\
\Cref{prop:caseI} & \lean{IsProfile.support_of_neg_aux}, \lean{IsProfile.support_of_neg}\\
Case II & \lean{IsProfile.support_of_neg_aux'}\\
\Cref{prop:clamp} & \lean{IsProfile.support_multiplier}\\
\eqref{eq:end-support} & \lean{IsProfile.support_one_zero}\\[4pt]
\multicolumn{2}{@{}l}{\emph{\Cref{sec:plt2,app:plt2}}}\\*[2pt]
\Cref{prop:low} & \lean{IsProfile.exists_burkholderMinorant_of_lt_two}\\
\Cref{lem:pone} & \lean{IsProfile.nonneg_of_one}\\
\Cref{lem:sublinear} & \lean{ConvexOn.eq_of_tendsto_div_zero}\\
\end{longtable}
}
